% Teil II: der durchgerechnete Weg.

% Gedankenschritt: Faktoren so wählen, dass beim Addieren eine Variable
% verschwindet; dann zurücksetzen und beide Gleichungen als Probe prüfen.

\subsection*{Gegeben und gesucht}
\begin{equation}
  -3x+4y=-1 \tag{I}
\end{equation}
\begin{equation}
  2x-5y=-4 \tag{II}
\end{equation}
gesucht: Lösungspaar $(x,y)$.

\subsection*{Ansatz und der entscheidende Gedankenschritt}
\textbf{Warum das Additionsverfahren?} Wir wollen $x$ loswerden. Die
$x$-Koeffizienten sind $-3$ und $2$; das kleinste gemeinsame Vielfache
ist $6$ (mit entgegengesetztem Vorzeichen). Multiplizieren wir (I) mit
$2\to -6x$ und (II) mit $3\to 6x$, so hebt sich beim Addieren $x$ auf:
$-6x+6x=0$. Es bleibt eine Gleichung nur in $y$.
Alternativ könnte man $y$ eliminieren (kgV der $y$-Koeffizienten $4$ und
$-5$ ist $20$); wir wählen die $x$-Variante.

\subsection*{Rechnung}
\begin{align}
  2\cdot(-3x+4y) &= -6x+8y=-2, \\[2pt]
  3\cdot(2x-5y) &= 6x-15y=-12 .
\end{align}
Addieren:
\begin{equation}
  -7y=-14 \;\Longrightarrow\; y=2 .
\end{equation}
In (II) einsetzen: $2x-5\cdot2=-4 \Rightarrow 2x-10=-4\Rightarrow 2x=6
\Rightarrow x=3$.

\textbf{Gedanke beim Rücksetzen:} Wir setzen $y=2$ in eine der
ursprünglich gegebenen Gleichungen (hier II), nicht in eine der
multiplizierten --- so bleibt die Rechnung direkt an der Ausgangslage
geankert.

\begin{ergebnis}
\begin{equation*}
  x=3,\qquad y=2 .
\end{equation*}
\end{ergebnis}

\begin{kontrolle}
\textbf{Probe in beiden Gleichungen:} (I): $-3\cdot3+4\cdot2=-9+8=-1$;
(II): $2\cdot3-5\cdot2=6-10=-4$. Beide Gleichungen erfüllt.

\textbf{Plausibilität:} Eingesetzt in die $-y$-Zeilengleichung
$-7y=-14$ stimmt $y=2$; die rechte Seite von (I), $-1$, wird mit dem
Paar $(3,2)$ exakt $-1$. Da der Schnittpunkt beide Geraden trifft, ist
das Paar eindeutig.

\textbf{Grenzfall/Idee:} Wäre die $x$-Elimination fehlgeschlagen (weil die
Gleichungen linear abhängig wären, also parallele Geraden), w\"are über das
Additionsverfahren $0=0$ (unendlich viele L\"osungen) bzw. $0=\text{Konst}$
(keine Lösung) herausgekommen. Dass wir einen eindeutigen Schnittpunkt
erhalten, bestätigt die Unabh\"angigkeit der beiden Geraden.
\end{kontrolle}

\begin{merksatz}
Beim Additionsverfahren systematisch eine Variable eliminieren (Faktoren
so, dass die Koeffizienten ein negatives Vielfachen-Paar bilden), lösen,
rücksetzen und die Lösung in \emph{beide} Gleichungen zur\"uckpr\"ufen.
\end{merksatz}
